% SHARD-ID: CA-31-GAUSSIAN-CURRENT-SYMBOL-EQUIVALENCE
% SHARD-TITLE: Gaussian Current-Symbol Equivalence in One Dimension
% SHARD-SUMMARY: Bridges the A4 real-space energy current to the CA-26 boost-time symbol for the nearest-neighbour scalar Gaussian chain.
% SHARD-SUMMARY: Checks only the translation-invariant integrated-current symbol, with the CA-29/CA-30 current orientation fixed explicitly.
% SHARD-SUMMARY: Leaves finite open-chain wave packets, dGamma(k), and continuum convergence as pending proof obligations.
% SHARD-KEYWORDS: Gaussian boson, energy current, boost-time symbol, Klein-Gordon, translation invariant, one dimension

\section{Gaussian Current-Symbol Equivalence in One Dimension}
\label{sec:gaussian-current-symbol-equivalence}

\subsection{Status, Sources, and Dependencies}

\textbf{Status: checked translation-invariant symbol bridge.}  This shard adds
the first bridge between the finite real-space current convention and the
momentum-symbol layer.  It is deliberately one-dimensional and
nearest-neighbour.  The checked statement is the equality of the
translation-invariant integrated-current symbol with the CA-26 boost-time
symbol for the scalar Klein-Gordon coefficients.  No finite open-chain momentum
operator, no finite periodic first moment, no \(\mathrm d\Gamma(k)\) theorem,
and no continuum convergence statement is claimed.

Dependencies: CA-24, CA-26, CA-29--CA-30, and CONVENTIONS.md (j), (k), (l).

% Convention: CONVENTIONS.md:196--239 -- finite periodic symbol checks do not
%   define periodic first moments or finite coordinate generators.
% Convention: CONVENTIONS.md:260--269 -- boost-time symbol
%   (1/2) partial_a omega(k)^2 and its sign convention.
% Convention: CONVENTIONS.md:379--384 -- open-chain current orientation
%   J_{j+1/2}=i[e_j,e_{j+1}] and continuity sign.
% Source: report/sections/26_gaussian_boson_generator_algebra.tex:133--150 --
%   boost-time residual and symbol target.
% Source: report/sections/30_gaussian_boson_1d_stress_candidates.tex:81--111 --
%   finite open-chain T_01 current candidate and continuity equation.
% Checked: src/GaussianBosonCurrents.jl:162--192 -- integrated current-symbol
%   helper and Klein-Gordon wrapper.
% Checked: test/runtests.jl:324--360 -- "Gaussian boson current-symbol
%   equivalence" invariant tests.

\subsection{Real-Space Current Orientation}

The finite open-chain A4 witness fixes the local current on an open bond by
\[
  J_{j+1/2}=i[e_j,e_{j+1}]
\]
with endpoint currents zero.  For a nearest-neighbour scalar chain with bond
coefficient \(b\), the checked closed form is
\[
  J_{j+1/2}
  =
  \frac b2\bigl(q_{j+1}p_j-q_jp_{j+1}\bigr).
\]
This sign is not adjustable in this shard: it is the CONVENTIONS.md (l)
orientation already used in the open-chain continuity test.

To compare with the boost-time symbol, pass from the local edge current to the
first-moment integrated current density.  In uniform spacing \(\epsilon\), the
nearest-neighbour first-moment coefficient contributes the spacing factor
\(\epsilon\).  The translation-invariant integrated current therefore has
symbol
\[
  \mathcal J_\epsilon(k)
  =
  -\epsilon b\sin(\epsilon k).
\]
This is a Fourier-symbol identity for the infinite or translation-invariant
nearest-neighbour expression.  It does not introduce a periodic coordinate
branch and does not construct an open finite-chain momentum operator.

\subsection{Klein-Gordon Equality}

For the one-dimensional nearest-neighbour lattice Klein-Gordon coefficients,
\[
  b=-\epsilon^{-2},
  \qquad
  \omega_{\epsilon,m}(k)^2
  =
  m^2+2\epsilon^{-2}(1-\cos(\epsilon k)).
\]
The integrated-current symbol becomes
\[
  \mathcal J_\epsilon(k)
  =
  \frac{\sin(\epsilon k)}{\epsilon}.
\]
The CA-26 boost-time symbol is
\[
  B_\epsilon(k)
  =
  \frac12\partial_k\omega_{\epsilon,m}(k)^2
  =
  \frac{\sin(\epsilon k)}{\epsilon}.
\]
Thus, at this checked symbol level,
\[
  \mathcal J_\epsilon(k)=B_\epsilon(k).
\]
The mass drops out because the current and the derivative of the onsite mass
term are both insensitive to \(m^2\).

\subsection{Julia Witness}

The helper
\[
  \texttt{nearest\_neighbor\_integrated\_energy\_current\_symbol}
\]
implements
\[
  -\epsilon b\sin(\epsilon k)
\]
for one-dimensional data only.  The Klein-Gordon wrapper inserts
\(b=-\epsilon^{-2}\).  The tests compare that wrapper against
\[
  \texttt{kg\_lattice\_boost\_time\_symbol}
\]
for several momenta and spacings.  They also include a sign sentinel:
for positive \(k\) and the Klein-Gordon negative bond, the current symbol is
positive and agrees with the boost-time target; reversing the bond sign gives a
negative symbol and fails the target comparison.  A small-spacing check compares
the current symbol with the centered continuum momentum representative \(k\)
near zero using the existing Gaussian small-spacing tolerance.

\subsection{Open Obligations}

This shard closes only the first symbol-level bridge.  The following remain
pending:
\[
\begin{array}{ll}
1.&\text{a finite open-chain wave-packet comparison with boundary residuals,}\\
2.&\text{an identification with the one-particle }\mathrm d\Gamma(k)
      \text{ momentum target,}\\
3.&\text{a periodic first-moment construction, which still needs a branch
      policy,}\\
4.&\text{a scaling or continuum convergence theorem.}
\end{array}
\]
Any compiler output citing this shard must therefore say ``current-symbol
equality'' rather than ``finite-chain momentum'' or ``continuum stress tensor''.
